\documentclass{article}
\usepackage{iclr,times}
\input{math.tex}
\usepackage{hyperref}
\usepackage{url}
\usepackage{booktabs}
\usepackage{multirow}
\usepackage{graphicx}
\usepackage{amsmath}
\usepackage{xcolor}

% Placeholder for a quantity that has NOT been measured yet. See NUMBERS.md: a number may be
% stated as fact only if it has a MEASURED row there. This macro exists so that an unmeasured
% slot is impossible to mistake for a result during editing.
\newcommand{\TODO}[1]{\textcolor{red}{\textbf{[TODO: #1]}}}

\title{One Canvas, Every Modality: Withholding Modalities \\ to Make Cross-Modal Completion Measurable}

\author{Anonymous}

\begin{document}
\maketitle

\begin{abstract}
Video-diffusion world models are increasingly asked to emit several modalities --- RGB, depth,
segmentation, surface normals, and robot actions --- from one checkpoint, with a prompt tag
selecting which. We show that in this design the prompt is not what selects the modality. Each
modality segment is handed a clean first latent frame, and that \emph{anchor} carries both the
modality's identity and its content: swapping the prompt tag moves accuracy by 1--2\%, while
removing the anchor costs roughly $9\times$. A model evaluated with anchors present is therefore
not demonstrating cross-modal generation, and it is being scored in a regime deployment cannot
supply, since a robot observes RGB only. We make the question answerable with two changes. First,
all five modalities for both camera views occupy a \emph{single} latent canvas --- ten segments,
110 latent frames --- so two modalities co-exist inside one attention operation for the first
time in this family of models. Second, we add a conditioning role that existing formulations
cannot express: a modality may be \emph{absent}, receiving no clean frame at all, so that it must
be produced from the others. A per-sample schedule over five regimes spends 45\% of samples with
the non-RGB modalities fully absent, which is simultaneously the deployment condition and the
only setting in which cross-modal completion is a prediction problem rather than a copy. We
report the cost of the canvas (2.5$\times$ sequence, 1.8$\times$ wall-clock per step, fitting in
48\,GB at four ranks), verify that the schedule takes effect in production from an independent
signature in step timing, and give a measurement protocol whose central object is a
\emph{mutual-completion matrix}: the quality of each modality as a function of how many of the
others are co-present. A flat matrix would show that the modalities merely share a decoder; only
a monotone one supports the claim that they supervise each other.
\end{abstract}

\section{Introduction}
\label{sec:intro}

A single generative model that emits RGB video, metric depth, semantic segmentation, surface
normals, and robot actions is attractive for embodied agents: one set of weights, one
representation, and auxiliary supervision that might regularise the action policy. The dominant
recipe in this family encodes every modality as an RGB image, packs a short sequence of
\emph{segments} into the latent canvas of a pretrained video diffusion transformer, and announces
in the text prompt which modalities the sequence contains. Switching the prompt switches the
task.

This paper starts from a measurement that undermines the usual reading of that recipe. In the
formulation we build on, every segment that is not fully supplied still receives its own first
latent frame as a clean condition --- an \emph{anchor}. We find that the anchor, not the prompt,
is what the model acts on. Substituting the prompt tag for a different modality's tag changes
accuracy by 1--2\%; deleting the anchor frames changes it by about $9\times$. Two consequences
follow, and both are structural rather than incidental:

\begin{enumerate}
\item \textbf{The capability is not what it appears to be.} If depth's own first frame is always
      supplied, ``the model generates depth'' and ``the model propagates depth's anchor forward''
      produce the same numbers. Cross-modal completion is untestable by construction.
\item \textbf{The evaluation regime is unreachable at deployment.} A robot at rollout time has
      RGB from its cameras and nothing else. There is no depth, segmentation, or normal first
      frame to hand over. A model scored with those anchors present is being asked, in the field,
      for precisely the $\sim\!9\times$-worse condition it never trained in.
\end{enumerate}

A third limitation is about the architecture rather than the masking. Because a sample draws one
template and a batch may not mix templates, \emph{two modalities never co-exist in a single
training step}. Such a model learns one set of weights taking turns at several jobs; it is never
given the opportunity to let depth inform RGB, or segmentation inform action, inside one forward
pass. Whatever mutual supervision the modalities could provide, this design cannot express.

We address all three with a deliberately simple construction, and then spend most of the paper on
how to measure whether it worked.

\paragraph{Contributions.}
\begin{itemize}
\item \textbf{A full-modality canvas.} All five modalities for both camera views are concatenated
      along the latent temporal axis into one sequence: ten segments, 110 latent frames, 28{,}160
      tokens (Sec.~\ref{sec:canvas}). Two modalities now co-exist in one attention operation. We
      report the realised cost: 2.5$\times$ the sequence, 1.8$\times$ the wall-clock per step,
      38.8\,GB peak at four ranks.
\item \textbf{An \emph{absent} conditioning role, and a schedule over it.} Beyond
      ``fully given'' and ``anchored'', a modality may receive no clean frame at all
      (Sec.~\ref{sec:faxis}). This is the role the anchor-dominance finding demands and the one
      deployment actually presents. We define five named regimes and a default mixture that
      spends 45\% of samples with the non-RGB modalities absent.
\item \textbf{A protocol whose headline is a falsifiable matrix.} The
      \emph{mutual-completion matrix} measures each modality's quality against the number of
      co-present sources (Sec.~\ref{sec:protocol}). Monotone means the modalities supervise each
      other; flat means they share a decoder. We state in advance which outcome would end this
      line of work.
\item \textbf{Two silent-failure modes in multimodal evaluation code}, both found here and both
      of a kind we expect to recur: a canvas layout transcribed in more than one place
      (Sec.~\ref{sec:pitfalls}). One mislabels which modality a metric belongs to; the other
      decodes a robot pose out of a segmentation map and reports it as a policy. Neither raises
      an error.
\end{itemize}

We are explicit about status. The canvas trains, the schedule demonstrably fires, and the
measurement machinery is built and unit-tested. The matrix itself, the closed-loop action
numbers, and the all-anchor control arm are \TODO{pending; see NUMBERS.md}. We report the
early, underpowered evidence we do have --- which currently says the anchor dependence is
\emph{not yet} broken --- rather than withholding it.

\section{Related work}
\label{sec:related}

\paragraph{Video diffusion as a world model.} Latent video diffusion transformers trained at
scale have become the default backbone for action-conditioned prediction. We build on the
Action~Images formulation, which renders a 6-DoF trajectory as three Gaussian heat maps in the
RGB channels and lets the same diffusion model that predicts video predict the action image,
recovering the pose by multi-view triangulation. This makes ``action'' just another segment,
which is what makes a single-canvas treatment of action and perception possible at all.

\paragraph{Multi-task vision foundation models.} Argus and related work show one backbone serving
many dense vision tasks, and also document the failure mode most relevant to us: beyond a data
budget, adding tasks costs accuracy on all of them. Our base configuration sits three orders of
magnitude below that budget, which is why we treat data scale as a first-class axis
(Sec.~\ref{sec:data}) rather than an afterthought.

\paragraph{Prompt-routed generation.} The idea that a text tag selects among generation tasks is
widespread. Our anchor-dominance measurement is a caution about how such claims are verified: if
each task's output is also given a clean first frame, the tag and the frame are confounded, and
the ablation that separates them (delete the anchor, keep the tag) is rarely run.

\paragraph{Privileged information.} Depth and ground-truth segmentation are privileged: available
in simulation, unavailable on a robot. The standard response is to distil a privileged teacher
into a deployable student. That framing presumes a gap to distil. Our construction instead puts
teacher and student in the same sequence and the same forward pass, and the matrix of
Sec.~\ref{sec:protocol} is what decides whether any transfer occurs.

\section{Preliminaries}
\label{sec:prelim}

\paragraph{Latent canvas.} The backbone is a 5B-parameter text-to-video diffusion transformer
(dim 3072, 30 blocks, 24 heads) over a causal video VAE with 48 latent channels, $16\times$
spatial and $4\times$ temporal compression. At $512^2$ and 41 frames one segment is a
$[48,11,32,32]$ latent; patchified with $(1,2,2)$ it contributes $11\times16\times16=2{,}816$
tokens.

\paragraph{Segments and roles.} A \emph{template} is a set of modalities
$\Pi\subseteq\{\textrm{video},\textrm{depth},\textrm{seg},\textrm{normal},\textrm{action}\}$. The
canvas is laid out view-major, modality-minor in a fixed canonical order, so
$\langle\textrm{depth}\rangle\langle\textrm{action}\rangle$ and
$\langle\textrm{action}\rangle\langle\textrm{depth}\rangle$ assemble identically and the model
need not spend capacity on permutation invariance. Training is masked diffusion: a boolean mask
marks \emph{condition} positions, those positions are restored to their clean latents after
noising, and the loss is taken only on the complement.

\paragraph{Action as pixels.} Action segments are not read from disk. The same 7-D trajectory is
reprojected through each view's own camera inside the forward pass, so the supervision target for
the action segment is pixel-identical across any template that contains it. This is what makes a
cross-modal consistency term well defined without alignment heuristics, and it is why the action
segment can be withheld and regenerated like any other.

\section{Method}
\label{sec:method}

\subsection{The full-modality canvas}
\label{sec:canvas}

We set $\Pi$ to all five modalities and keep both views, giving ten segments:
\begin{center}
\texttt{V0 D0 S0 N0 A0 | V1 D1 S1 N1 A1}
\end{center}
Concatenation is along the latent temporal axis --- the same axis the backbone's rotary position
embedding indexes. The model therefore does not see ``modalities''; it sees one 110-frame video
and must infer from position and content which stretch is which. Nothing in the architecture is
modality-aware, and no modality embedding is added.

This is the smallest change that makes mutual supervision expressible: every pair of modalities
is now inside the same dense self-attention.

\paragraph{Cost.} The canvas is 28{,}160 tokens against 11{,}264 for the four-segment templates:
$2.5\times$ the sequence and, because attention is quadratic, $3.6\times$ the transformer FLOPs.
Measured throughput is 36--37\,s/step against 20.5\,s/step, i.e.\ $1.8\times$ --- less than the
FLOP ratio, because a large constant dominates. Solving the step-time pair between the full
canvas (110 frames, 37\,s) and the collapsed policy canvas (30 frames, 20\,s) under the
transformer's own scaling gives a fixed overhead of $\approx\!16.5$\,s/step, 45\% of a step. That
constant is VAE encoding (fourteen encodes per step) plus the CPU-offloaded optimiser, not
attention. It is the dominant remaining inefficiency and is addressed in
Sec.~\ref{sec:limitations}.

Peak memory is 38.8\,GB of 48\,GB at four data-parallel ranks, measured over 2{,}400
one-second samples with no sample above 43\,GB. Gradient sharding is what buys this: at two ranks
the same canvas is projected to exceed the card.

\subsection{Three conditioning roles}
\label{sec:roles}

Each segment takes exactly one role:
\begin{description}
\item[\textsc{given}] every latent frame is a clean condition; the modality is supplied.
\item[\textsc{anchored}] only frame 0 is clean; frames $1..T$ are predicted. This is the
      historical default for every non-supplied segment.
\item[\textsc{absent}] \emph{no} frame is clean. The segment is noise throughout and must be
      produced from the rest of the sequence.
\end{description}
\textsc{absent} is new. Its absence from prior formulations is exactly why anchor dominance went
unmeasured: with only the first two roles, no sample ever asks a modality to be generated without
its own first frame.

Two constraints are load-bearing. Roles are assigned \emph{per modality and applied to both
views}: a per-view assignment would let view~0's depth anchor view~1's depth, making ``absent'' a
fiction. And at least one segment must be non-absent, or the canvas has no clean pixel anywhere
and the sample degenerates to text-to-video; we assert this rather than rely on the loss to
notice.

\subsection{The modality-dropout schedule}
\label{sec:faxis}

Each sample draws one of five regimes, which assign the roles above:

\begin{table}[t]
\centering
\small
\begin{tabular}{llccc r r}
\toprule
Regime & $p$ & \textsc{given} & \textsc{anchored} & \textsc{absent} & frames & cond.\ frames \\
\midrule
\texttt{full\_anchor} & 0.30 & --- & all five & --- & 110 & 10 \\
\texttt{rgb\_only}    & 0.30 & --- & video & d,s,n,a & 110 & 2 \\
\texttt{rgb\_given}   & 0.15 & video & --- & d,s,n,a & 110 & 22 \\
\texttt{one\_out}     & 0.15 & --- & four & one, drawn uniformly & 110 & 8 \\
\texttt{policy}       & 0.10 & --- & video (1 frame) & d,s,n (1 frame) & 30 & 4 \\
\bottomrule
\end{tabular}
\caption{The default schedule (\textsc{f1}). \texttt{rgb\_only} is the deployment regime. The
draw in \texttt{one\_out} includes \emph{video} and \emph{action}, so depth$\rightarrow$RGB and
observations$\rightarrow$action also receive gradient; restricting it to the auxiliary modalities
would train only ``RGB explains everything'' and the word \emph{mutual} would not be earned.
\texttt{policy} reproduces the canvas a closed-loop replan actually presents. Condition-frame
counts are measured, not nominal.}
\label{tab:faxis}
\end{table}

The realised mixture over 4{,}000 draws is $(0.307, 0.316, 0.139, 0.138, 0.100)$, matching the
nominal values within sampling noise. The schedule \emph{replaces} rather than composes with the
prior conditioning axes: composing them would allow a modality to be simultaneously the supplied
source and absent.

\paragraph{The prompt names the canvas, not the supply.} The text prefix lists all five
modalities under every regime. Given that the tag was measured to carry little the model acts on,
encoding absence in it would be decoration; worse, it would make the text distribution depend on
a draw that any evaluator must then reproduce exactly.

\paragraph{Action dilution, measured.} A four-segment template puts 50\% of predicted latent
positions on action; this canvas puts 26.3\% (0.20 in the anchored regimes, 0.77 in
\texttt{policy}, which concentrates its loss there). The naive estimate is 20\%, i.e.\ a
$2.5\times$ dilution; the realised figure is $1.9\times$. We accept this and state it as a
confound rather than introducing per-modality loss weights, which would add a free parameter and
make the arm incomparable to every other.

\section{Measurement protocol}
\label{sec:protocol}

\subsection{The mutual-completion matrix}

The object that decides whether this direction is worth pursuing is, for each modality $X$, its
quality against the number of co-present sources:
\begin{center}
\begin{tabular}{lll}
\toprule
Setting & $X$'s role & co-present sources \\
\midrule
\texttt{full\_anchor} & \textsc{anchored} & --- ($X$ has its own frame; upper bound) \\
\texttt{one\_out}:$X$ & \textsc{absent} & four \\
\texttt{rgb\_only}    & \textsc{absent} & one (video) \\
\bottomrule
\end{tabular}
\end{center}
If $X$ improves monotonically with co-present sources, the modalities are informing each other.
\textbf{If the matrix is flat, they are merely sharing a decoder}, the premise of this direction
fails, and the correct action is to stop rather than scale --- the same verdict a prior
privileged-distillation plan in this codebase reached when its assumed teacher--student gap
turned out not to exist. We commit to that reading in advance.

Surface normals are the sharpest probe of the five. The normal encoder is a deterministic
function of depth and focal length, so when normals are absent and depth is present there is a
\emph{correct} answer to recover, and failure cannot be attributed to ambiguity.

\subsection{What to report for each modality}

Depth is scored by absolute relative error \emph{together with decode coverage}, the fraction of
pixels still on the depth codec's colour path. This pairing is not optional: a generated image
that leaves the path decodes to \texttt{NaN}, those pixels are excluded, and the error is then
computed over whatever remains. Coverage is the difference between ``the depth is inaccurate'' and
``the output is not a depth image''. Segmentation is scored by macro-mIoU plus the count of
\emph{unknown} pixels, and normals by mean cosine. For action we report both an open-loop
position error and closed-loop task success, since the former saturates in ways the latter does
not.

\subsection{Pitfalls in multimodal evaluation code}
\label{sec:pitfalls}

The canvas layout is an implicit contract that many components depend on, and we found it
transcribed independently in two evaluation paths. Both failures are silent.

\begin{enumerate}
\item A scorer's span table carried the comment ``mirrors the pipeline exactly''. It went stale
      the moment \textsc{absent} was added, and a stale span table does not crash --- it reports
      one modality's metric under another modality's name.
\item A closed-loop decoder assumed four equal segments and sliced the second and fourth as the
      action heat maps. On the ten-segment canvas those slices are \emph{segmentation}: the
      decoder would have triangulated a 6-DoF pose out of a segmentation map and reported it as a
      policy. Both are RGB images, so nothing raises.
\end{enumerate}
The fix is not to re-synchronise the copies but to delete them: one function maps
(template, regime) to roles, and the pipeline and every scorer call it. We verified that the
four-segment spans are unchanged and that the ten-segment action spans move from indices
$\{1,3\}$ to $\{4,9\}$. We recommend treating any ``mirrors X exactly'' comment as a latent
defect.

\section{Experiments}
\label{sec:exp}

\subsection{Setup}
\label{sec:data}

\paragraph{Data.} RLBench, 16 tasks, four calibrated cameras at $512^2$, Colosseum-style
randomisation of camera pose, table colour and texture, background texture and light colour. Two
trees differing \emph{only} in scale: 788 and 3{,}960 variation-0 training episodes (5.0$\times$),
each with 240 held-out episodes from unseen variations. Depth comes from the simulator, normals
analytically from depth, and segmentation from the simulator's object-handle map resolved to
functional roles.

Scaling the tree exposed an annotation gap invisible at the smaller scale. The room ceiling
appears in the handle map of 3 of 4{,}200 episodes --- never in the 788-episode tree --- because
camera randomisation occasionally draws an elevation that sees past the wall tops. A union over
each task's handles, itself present to fix a different sampling gap, then propagated those three
episodes to all 540 episodes of the affected tasks, and the startup check that refuses
\emph{unknown} pixels would have rejected the whole tree. We report it because it is a
general hazard of data scaling, and because verifying the fix required showing the smaller tree
was bit-identically unaffected, which we did by exhaustive scan before changing anything.

\paragraph{Verification.} Every episode was checked structurally; 99 episodes (396
episode--view pairs) were decoded through the repository's own codecs. Depth round-trip AbsRel is
$7\times10^{-4}$ (median), normal cosine $1.0000$, \emph{unknown} pixels zero. Running the
identical script on the smaller tree yields the same statistics, establishing that the two trees
differ in scale alone.

\paragraph{Training.} Warm-started from a released Action~Images checkpoint, four data-parallel
ranks, global batch 4, 5{,}000 steps (20{,}000 samples), $512^2$, 41 frames, lr $5\times10^{-7}$,
ZeRO-2 with CPU-offloaded optimiser, gradient checkpointing.

\subsection{Does the canvas train, and does the schedule fire?}

Both arms reached 5{,}000 steps without instability. Independently of any logged metric, the
schedule is visible in step timing: \texttt{policy} collapses the canvas from 110 to 30 latent
frames, and 7--8\% of steps take 20--21\,s against 37\,s for the rest --- matching the nominal
10\% within noise. We note this because a mask schedule that silently fails to take effect
produces a normal-looking loss curve, and the repository's own guidance is that configuration
changes of this kind must be confirmed from something other than the configuration file.

\subsection{Per-modality quality under each conditioning regime}

\begin{table}[t]
\centering
\small
\begin{tabular}{lrrrrr}
\toprule
Regime & video PSNR & depth AbsRel & depth coverage & normal cos & seg unknown px \\
\midrule
\texttt{full\_anchor} & 15.38 & \textbf{0.121} & \textbf{0.9995} & \textbf{0.832} & 35{,}119 \\
\texttt{rgb\_only}    & 14.83 & 0.395 & 0.409 & $-0.286$ & 4{,}218 \\
\texttt{rgb\_given}   & 30.71$^\dagger$ & 1.202 & 0.043 & $-0.345$ & 10{,}371 \\
\texttt{policy}       & 32.30 & 0.844 & 0.9992 & 0.718 & \textbf{0} \\
\bottomrule
\end{tabular}
\caption{Base-tree arm at step 2{,}000 of 5{,}000, \textbf{one episode, one seed}. A direction
indicator, not a result. $^\dagger$video is \textsc{given} here, so its PSNR is a VAE round-trip
floor, not generation quality. The scaled-tree arm at 5{,}000 steps is \TODO{eval queued}.}
\label{tab:regimes}
\end{table}

Table~\ref{tab:regimes} is early and underpowered, and we read only its largest effect. Depth
coverage falls from $0.9995$ when depth is \textsc{anchored} to $0.041$--$0.409$ when it is
\textsc{absent}, and normal cosine becomes \emph{negative}, i.e.\ the predicted normals point the
wrong way --- worse than the $\approx 0$ of an uninformative prediction. The apparent paradox that
\texttt{rgb\_given} (more information supplied) scores worse on depth than \texttt{rgb\_only} is
resolved by coverage: its AbsRel is computed over 4.3\% of pixels, the rest having left the codec
path. The honest summary is that \textbf{at 40\% of training the anchor still carries essentially
all of the ``which modality am I'' signal}, which is the dependence the schedule exists to break
and has not yet broken.

One contrary observation is worth recording. Under \texttt{policy}, depth is also
\textsc{absent}, yet coverage is $0.9992$, normal cosine $0.718$, and \emph{unknown} pixels zero.
The difference is span: that canvas is 30 frames and the depth frame is adjacent to the video
frame, whereas \texttt{rgb\_only} must generate eleven depth frames at temporal positions
11--21. Short-span anchor-free generation works; long-span does not yet. This is consistent with
either undertraining or rotary-position extrapolation --- the tenth segment sits at positions
99--109, while the backbone was pretrained to $\approx\!31$ latent frames and the warm start saw
44 --- and a per-segment loss breakdown is logged to separate them \TODO{not yet extracted}.

\subsection{Mutual completion}
\TODO{the matrix of Sec.~\ref{sec:protocol}: one-out sweep over all five modalities, on the
scaled-tree arm at 5{,}000 steps. Queued; every GPU on the host is currently held by other
tenants.}

\subsection{Closed-loop action}
\TODO{success rate over three tasks $\times$ 10 trials, asked through the ten-segment canvas
under \texttt{rgb\_only}, i.e.\ one RGB frame per view and nothing else. Queued.}

\subsection{Ablations}
\TODO{(i) the \textsc{f0} all-anchor control, which isolates what the schedule itself buys ---
same template, no absent role; (ii) the 788 vs 3{,}960 episode comparison at matched exposure;
(iii) normals removed, which tests whether an analytically redundant modality earns its 20\% of
the sequence.}

\section{Limitations}
\label{sec:limitations}

\paragraph{The central claim is not yet established.} What we can state is that the canvas
trains, the schedule fires, the protocol is implemented and unit-tested, and the early evidence
says anchor dominance survives 2{,}000 steps. Whether it breaks with training, with data, or at
all is open.

\paragraph{Data scale.} Even the scaled tree is 3{,}960 episodes of 16 tasks from one simulator.
The multi-task literature's own finding is that adding tasks beyond a data budget costs accuracy
on all of them, and we are far below that budget. A negative matrix at this scale would not
license a claim about the construction in general.

\paragraph{Statistical power.} Every fusion number here is one episode and one seed. The
repository's history includes a headline number that reversed under a properly powered re-test,
and we treat our own Table~\ref{tab:regimes} with that precedent in mind.

\paragraph{Efficiency.} 45\% of each step is VAE encoding and the CPU-offloaded optimiser, not
the transformer. Every stream is encoded at full length before the mask is known, so the short
\texttt{policy} canvas saves transformer time only. A pre-computed latent cache would remove most
of this; the obstacle is that the sampled temporal window is drawn per item, so the cache must be
keyed on the window or the window must be pinned.

\paragraph{Position encoding.} Ten segments place the last one at rotary positions 99--109, well
beyond both the backbone's pretraining range and the warm start's. We deliberately did not change
the position encoding, to keep the arm warm-start compatible; a per-segment reset plus a learned
modality embedding is the natural remedy and belongs to a from-scratch run.

\section{Conclusion}

Prompt-routed multimodal generation has been read as the prompt selecting a capability. We find
the supplied anchor frame doing that work instead, which makes cross-modal completion both
unverifiable and untrained in the only regime deployment offers. Putting every modality in one
latent canvas makes mutual supervision expressible, and adding an \emph{absent} conditioning role
makes it measurable. We have built and verified both, priced the canvas, and specified a
falsifiable matrix together with the outcome that would end the line of inquiry. The matrix
itself is the next measurement, not a claim we make here.

\bibliographystyle{iclr}
\bibliography{reference}

\appendix
\section{Canvas layouts per regime}
Measured, not nominal; \texttt{F}~= given, \texttt{I}~= anchored, \texttt{-}~= absent,
\texttt{1}~= segment collapsed to a single latent frame.
\begin{center}
\small
\begin{tabular}{ll}
\toprule
\texttt{full\_anchor} & \texttt{V0:I D0:I S0:I N0:I A0:I | V1:I D1:I S1:I N1:I A1:I} \\
\texttt{rgb\_only}    & \texttt{V0:I D0:- S0:- N0:- A0:- | V1:I D1:- S1:- N1:- A1:-} \\
\texttt{rgb\_given}   & \texttt{V0:F D0:- S0:- N0:- A0:- | V1:F D1:- S1:- N1:- A1:-} \\
\texttt{one\_out}:d   & \texttt{V0:I D0:- S0:I N0:I A0:I | V1:I D1:- S1:I N1:I A1:I} \\
\texttt{policy}       & \texttt{V0:I1 D0:-1 S0:-1 N0:-1 A0:I | V1:I1 D1:-1 S1:-1 N1:-1 A1:I} \\
\bottomrule
\end{tabular}
\end{center}

\section{Reproducibility}
The canvas is selected by the template string
\texttt{video+depth+segmentation+normal+action}; the schedule by a single flag with presets
\textsc{f0}/\textsc{f1}/\textsc{f2}. Role assignment for both training and evaluation comes from
one function, so an evaluation can only request a conditioning the training schedule can express.
Diagnostics logged per step include per-segment loss normalised by the step's own loss (absolute
per-segment values are dominated by the timestep draw and are not comparable across steps), a
one-hot over the drawn regime, and the allocator's peak memory. Periodic \texttt{nvidia-smi}
polling is insufficient for the last of these: a 4\,s poll missed a transient by 4.5\,GB in our
own measurements.

\end{document}
